\chapter{Interrupts and Reset}
\label{ch:interrupts}

The CPU has five interrupt inputs: four maskable lines, \activelow{IRQ1} to \activelow{IRQ4}, and the non-maskable
\activelow{NMI}. It has a reset input and a reset output. This chapter says what drives each line and how reset reaches
every chip.

\section{Interrupt Lines}
\label{sec:interrupt-lines}

\begin{table}[H]
	\centering
	\small
	\begin{tabular}{llllp{4.6cm}}
		\toprule
		\textbf{Line}    & \textbf{Level} & \textbf{Source}   & \textbf{Condition}     & \textbf{Behaviour}                                                  \\
		\midrule
		\activelow{NMI}  & 5              & SIRCADA           & Vertical sync          & One-cycle pulse at the start of vertical blanking; edge-triggered   \\
		\activelow{IRQ4} & 4              & SIRCADA           & Line compare           & Level-sensitive; held until the program acknowledges it             \\
		\activelow{IRQ3} & 3              & SIRCAS            & A sequence is complete & Level-sensitive; held until the program clears the completion flags \\
		\activelow{IRQ2} & 2              & Serial device     & Character or timer     & Level-sensitive; held until the program clears the cause            \\
		\activelow{IRQ1} & 1              & Not yet specified &                        &                                                                     \\
		\bottomrule
	\end{tabular}
	\caption{Interrupt Assignments}
	\label{tab:interrupt-lines}
\end{table}

Every source is an open-drain output, so a line may be shared by several devices and is released by all of them when
none asks for service. The pull-up that holds each line high is part of the system. Its value is not yet specified. A
handler on a shared line reads the status of each device to find which one is asking. The display adapter has a bit for
this purpose, \texttt{LIP}, in its Status register (\videomanual{}, timing chapter).

The CPU samples the lines on every rising edge of \texttt{CLKI} except during reset, and a level-sensitive line must
stay asserted until the handler has acknowledged it (\hwmanual{}, signal descriptions chapter). A device that withdraws
a request before it is latched is not recognised. The two devices follow this rule: the display adapter holds
\activelow{IRQ4} until the program writes the \texttt{LINE} register, and SIRCAS holds \activelow{IRQ3} until the
program writes ones to \texttt{SEQ\_DONE}. The serial device follows it too, as the next section describes.

\section{The Serial Device's Interrupt}
\label{sec:serial-interrupt}

The serial device has one interrupt output, \activelow{IRQ}, wired to the CPU's \activelow{IRQ2}. Eight conditions can
drive it, and a mask register chooses which do. They are, for each channel, a received character (or a full receive
buffer), room in the transmit buffer, and a change in a break; and for the device, the counter/timer reaching zero and
a change on the input port. The interrupt status register shows which are true, so one handler for \activelow{IRQ2}
reads that register and serves each condition in turn. The controller input pins are polled and raise no interrupt.

\activelow{IRQ} is level-sensitive and stays asserted until the program clears the cause: it reads the received
character, writes a character to send, issues the reset command for a break change, stops the timer, reads the input port
change register, or masks the condition. It is negated within 300 ns of the end of that access. This is what the CPU needs of a
level-sensitive line. The CPU has no interrupt-acknowledge cycle, so the DUART's acknowledge input is tied high and its
vector register is not used. The handler is the one at the \activelow{IRQ2} entry of the vector table, words \texttt{0x0080} and
\texttt{0x0081} (\archmanual{}, exceptions chapter).

\subsection{Waiting Behind Higher Levels}

\activelow{IRQ2} ranks below the vertical sync, line-compare and SIRCAS interrupts, so a request waits while one of their handlers
runs. The receiver holds four characters. At 38,400 baud a character takes 260~$\mu$s, so a request can wait about a
millisecond before characters are lost. Vertical blanking is 2.41~ms long in an NTSC system, so a vertical sync handler that
uses all of it can cause an overrun at that speed. A program keeps those handlers short, or uses a lower speed: at 19,200 baud
the wait allowed is about 2~ms, and at 9,600 about 4~ms. The transmitter takes two characters, and needs service only to keep
the line busy.

The counter/timer is a 16-bit counter that counts down from the 3.6864 MHz clock or a division of it and interrupts
when it reaches zero. It makes a timer interrupt on \activelow{IRQ2} that does not depend on the picture, for time-outs
on the link, for example.

\section{Vertical Sync Is Special}

The vertical sync pulse of the display adapter is a single \texttt{CLKI} cycle on the non-maskable line, so it needs no
acknowledgement and cannot be masked. The CPU does not miss it because it samples the line on every edge. A pulse that
arrives while its own handler is still running raises the Level Five Hardware Exception Conflict fault (\archmanual{},
exceptions chapter), so the handler must return within a frame: 16.64~ms in an NTSC system, 20.00~ms in a PAL system. A
DMA transfer cannot be interrupted, and so it delays the handler if it is running at the time (\videomanual{}, timing
chapter).

\section{Reset}
\label{sec:reset}

Reset has one input and one output at the CPU, and each of the two other chips has a reset input:

\begin{table}[H]
	\centering
	\begin{tabular}{lll}
		\toprule
		\textbf{Net}     & \textbf{Driven by}             & \textbf{Received by}                                              \\
		\midrule
		\activelow{RSTI} & The power-on and reset circuit & CPU \activelow{RSTI}                                              \\
		\activelow{RSTO} & CPU                            & SIRCADA \activelow{RST}, SIRCAS \activelow{RST}, the boot overlay \\
		\bottomrule
	\end{tabular}
	\caption{The Reset Nets}
	\label{tab:reset-nets}
\end{table}

The reset circuit, which asserts \activelow{RSTI} at power-up and when the console is reset by hand, is not yet
specified.

The CPU asserts \activelow{RSTO} while \activelow{RSTI} is asserted, and for six \texttt{CLKI} cycles after it is
released, before it fetches the reset vector. It asserts it again for the same six cycles after a software reset (the
\mnemonic{RSET} instruction), after a failed double-fault vector fetch, and before it retries a failed reset-vector
fetch (\hwmanual{}, signal descriptions chapter). \activelow{RSTO} is wired to the reset inputs of both other chips and
sets the boot overlay (Section~\ref{sec:boot-overlay}), so a reset of the CPU is a reset of the whole console and
leaves the system ROM answering at the low words. A program that resets the CPU must expect to start again with the
display and the sound in their reset states, and must set both up again.

What reset does to each chip is in its own manual. In short: SIRCAS goes silent, its registers take their reset values
and the contents of the audio memory are lost (\audiomanual{}, registers chapter); SIRCADA comes up blank, in the mode
it has at reset, with every writable register zero but a few (\videomanual{}, registers chapter). A chip needs its
reset input held for a minimum time. SIRCADA needs six cycles (\videomanual{}, signal timing chapter), which the CPU's
six-cycle hold provides. The minimum for SIRCAS is not yet specified.
